\section{Men without Qualities}

\begin{quotex}
But because thou art lukewarm, and neither cold, nor hot, I will begin to vomit thee out of my mouth.
\flright{\textsc{Revelation 3:16}}
\end{quotex}


\textbf{Robert Musil}`s novel, \emph{A Man without Qualities}, is about


\begin{quotex}
a 32-year-old mathematician named Ulrich who is in search of a sense of life and reality but fails to find it. His ambivalence towards morals and indifference to life has brought him to the state of being ``a man without qualities'', depending on the outer world to form his character. A kind of keenly analytical passivity is his most typical attitude.
\flright{\textit{The Man Without Qualities}\footnote{\url{https://en.wikipedia.org/wiki/The_Man_Without_Qualities}}}
\end{quotex}


The man without qualities suffers from acedia (state of listlessness, indifference to one's position or condition in the world) and abulia (lack of will). There is no lack of intelligence, yet his intellectual strivings have no purpose so can neither form a plan of action nor lead to any peace or certainty. Here is a self-description of such a man, which recently passed through my news feed. He identifies as an

\begin{quotex}
agnostic with pagan sympathies
\end{quotex}


After a lifetime of thinking, that is a rather meagre result. There is no commitment to anything: neither to a purpose to the world nor to an acknowledgment of chaos. What is a pagan sympathy? If a man is squeamish about human sacrifice or infanticide, how sympathetic can he be to the pagans? The myths and legends of the ancient pagans are replete with patricide, matricide, squabbles, deceit, treachery, and so on. The same goes for the stories in the Edda. They are hardly sympathetic.

\paragraph{A Misnomer}


In an online public discussion, someone asserted:

\begin{quotex}
Traditionalism posits a return to a mentality through intellectual means to which we, as modern people, cannot return, any more than a 60-year-old can return to the mentality of a 6-year-old. (Which is not to denigrate either mentality.)
\end{quotex}


It's a sad sign of the times that such a sentence is considered a marker of intelligence in some circles. It would be forgivable except that it was authored by someone who once claimed to be ``sympathetic'' toward tradition. We dealt with the misnomer ``traditionalism'' some 9 years ago in Tradition and Traditionalism\footnote{\url{https://gornahoor.net/?p=383}} and in The Absurdity of Traditionalism\footnote{\url{https://gornahoor.net/?p=728}} \textbf{Rene Guenon} addresses this topic nearly a century ago:

\begin{quotex}
People described as Traditionalists ... only have a sort of tendency or aspiration towards tradition \emph{without really knowing anything at all about it}; this is the measure of the distance dividing the ``traditionalist'' spirit from the truly traditional spirit,
\emph{for the latter implies a real knowledge}, and indeed in a sense it is the same as that knowledge.
\flright{\textsc{Rene Guenon}, \emph{The Reign of Quantity and the Signs of the Times}}
\end{quotex}


``Tradition'' must be understood in its literal sense as something ``handed down'' or ``given across'', not as the ``repetition of the past'' for its own sake. What is handed down, then, is a transcendent knowledge which cannot be considered as just another perspective or worldview. Hence, the ``modern'' mentality must certainly be denigrated since it is the opposite of authentic knowledge. Rather it is a symptom of intellectual confusion. Guenon continues:

\begin{quotex}
words are applied to things which they do not fit in any way, and sometimes in a sense directly opposed to their normal meaning. This is one of the most obvious symptoms of the intellectual confusion which reigns everywhere in the present world; but it must not be forgotten that this very confusion is willed by that which lies hidden behind the whole modern deviation.
\end{quotex}


Harsh words from the Master. The man without qualities may believe he is directing his life, but, as Guenon observes, his vaunted modern mentality is formed by hidden spiritual forces beyond his awareness. The modern mind has sunk into materiality, making it unable to perceive anything higher.

\paragraph{Archaic Consciousness and Eternal Return}


The correct formulation of a question often leads to the correct answer. The question, then, should be rephrased as whether a return to archaic consciousness, in \textbf{Mircea Eliade}'s sense of the term, is possible. He explains:

\begin{quotex}
In \emph{imitating} the exemplary acts of a god or of a mythical hero, or simply by recounting their adventures, the man of an archaic society detaches himself from profane time and magically re-enters the Great Time, the sacred time.
\flright{\emph{Myths, Dreams and Mysteries}}
\end{quotex}


Of course, the modern man, unable to take such myths at face value, does so in a different manor from the primitive man. Paul Ricoeur calls this re-enactment the second naïveté. In any case, knowledge of sensory space-time events can never be more than a matter of opinion (\emph{doxa}). Hence, sacred time refers to events in consciousness.
\textbf{Hans Leisegang} makes this clear:


\begin{quotex}
Every myth expresses, in a form narrated for a particular case, an eternal idea, which will be intuitively recognised by him who re-experiences the content of the myth.
\flright{\textsc{Die Gnosis}}
\end{quotex}


\textbf{Valentin Tomberg} in \emph{Meditations on the Tarot} points out the relationship between myths and archetypes:

\begin{quotex}
myths reveal the archetypes which manifest themselves endlessly in history and in each individual biography -- they are mythological symbols pertaining to the domain of time.
\end{quotex}


There is a method of Hermetic Meditation\footnote{\url{https://gornahoor.net/?p=898}} that enables one to re-enter into that mythological or sacred time. The goal is to regain the state of consciousness that existed before the Fall. Eliade provides this example:

\begin{quotex}
The most representative mystical experience of the archaic societies, that of shamanism, betrays the \emph{Nostalgia for Paradise}, the desire to recover the state of freedom and beatitude before `the Fall'.
\end{quotex}


This is true not only of shamanism, but also of the Hermetic teachings.

\paragraph{Intellectual Conversion}


Eliade offers this explanation of the Traditional point of view:

\begin{quotex} 
things acquire their reality, their identity, only to the extent of their participation in a transcendent reality.
\flright{\textit{Myth of the Eternal Return}}
\end{quotex}


Clearly, this transcendent reality is the realm of ideas and Plato is the philosopher who made this explicit. Hence, Eliade continues:

\begin{quotex}
Plato could be regarded as the outstanding philosopher of `primitive mentality,' that is, as the thinker who succeeded in giving philosophic currency and validity to the modes of life and behaviour of archaic humanity.
\end{quotex}


\textbf{Julius Evola} also distinguishes between the worlds of being and becoming. In the first chapter of \emph{Revolt Against the Modern World}, this is made clear and explicit:

\begin{quotex}
In order to understand both the spirit of Tradition and its antithesis, modern civilization, it is necessary to begin with the fundamental doctrine of the two natures. According to this doctrine there is a physical order of things and a metaphysical one; there is a mortal nature and an immortal one; \emph{there is the superior realm of being and the inferior realm of becoming}. Generally speaking, there is a visible and tangible dimension and, prior to and beyond it, an invisible and intangible dimension that is the support, source, and true life of the former.
\end{quotex}


The modern mentality is tied to the physical order of things, or the inferior realm of becoming. He is the cave dweller in Plato's myth of the cave, and can only see the shadow of reality on the wall. By assiduously wrestling with Plato's text or similar writings, the modern man may undergo an intellectual conversion. He then escapes from the cave and discovers the Sun. Sadly, few are willing to make such efforts or may actually be incapable of it. It is the nature of error not to recognize itself as error. Such a man will disbelieve any rumours about the Sun and is like those living in The Country of the Blind\footnote{\url{https://gornahoor.net/?p=7992}}.

Hence, Guenon opined that the future would be open to a small remnant, not by the restoration of past glory, but rather by adapting those principles to the current situation. He elaborates in \emph{The Crisis of the Modern World}:

\begin{quotex}
If it be objected that Christianity itself, in our time, is no longer understood in its profound meaning, we should reply that it has at least kept in its very form all that is needed to provide the foundation of which we have been speaking. The least fantastic venture, in fact the only one that does not come up against immediate impossibilities, would therefore be an attempt to restore something comparable to what existed in the Middle Ages, with the differences demanded by modifications in the circumstances; and for all that has been completely lost in the West, it would be necessary to draw upon the traditions that have been preserved in their entirety, as we stated above, and, having done so, to undertake the task of adaptation, which could be the work only of a powerfully established intellectual elite.
\end{quotex}

\paragraph{Postscript} The point of this is not to single out any one person, but rather to answer a legitimate question in a complete way. There many so called `radical traditionalists' or those `influenced by tradition', etc., without any real understanding. Just as bad money drives out the good, so too do bad ideas drive out the truth.

We have defined the ideas of Tradition, the modern mentality, and archaic consciousness with precision. Ultimately, the answer to the original question cannot be proven by logic, but only by personal experience. The alternative is to live with ennui, confusion, and purposelessness.

\flrightit{Posted in 18-02-2019 by Cologero}